\subsection{Unipotent residual image leaves only the order-twelve versus order-six pair}

\begin{lemma}[The two unit depths over absolute ramification two]
\label{lem:B-RD39-units}
\status{Proved. Complete root and independent proof, original-source, and exact-certificate audits passed.}
Let $K$ be complete of mixed characteristic $(0,3)$, with algebraically
closed residue field, and normalize $v_K(3)=2$. Choose $\pi^2=3$.
Every unit in $1+\pi^3\mathcal O_K$ is a cube. A noncube unit $D$
therefore has a well-defined cube-approximation depth $s\in\{1,2\}$:
for every unit $t$ with $v_K(t^3-D)>0$, one has
$v_K(t^3-D)=s$. The inverse unit $D^{-1}$ has the same depth.

If $L/K$ is totally ramified of degree three, and a cube
$u\in(L^\times)^3$ satisfies $0<v_L(u-1)<9$, then
$v_L(u-1)$ is either three or six.
\end{lemma}
\begin{proof}
For $b\in\mathcal O_K$, the equation
\[
       (1+\pi w)^3=1+\pi^3b
       \quad\Longleftrightarrow\quad
                     w^3+w+\pi w^2=b
                                                        \tag{B-RD39.1}
\]
has a solution in $\mathcal O_K$: its reduction has a root over the
algebraically closed residue field and its derivative is a unit.
This proves the first assertion. Choose a unit $t$ whose residue is
the unique cube root of $\overline D$. The positive valuation of
$t^3-D$ must be one or two, since valuation at least three would
make $D/t^3$, and hence $D$, a cube. Any two such $t$ differ by a
multiple of $\pi$, and their cubes differ by a multiple of $\pi^3$.
The positive valuation is therefore independent of $t$.
Replacing $t$ and $D$ by their inverses preserves it.

For the last assertion write $u=x^3$. Its residue is one, so
$x=1+w$ with $r=v_L(w)\ge1$. Since $v_L(3)=6$, the three terms in
$u-1=3w+3w^2+w^3$ have valuations $6+r,6+2r,3r$.
For $r=1,2$ their minimum is uniquely $3r$. For $r\ge3$ every term
has valuation at least nine. This proves the claim.
\end{proof}

\begin{proposition}[B-RD39: the cyclic-three residual case]
\label{prop:B-RD39}
\status{Proved. Complete root and independent proof, original-source, and exact-certificate audits passed.}
Let $K$ be as in Lemma~\ref{lem:B-RD39-units}. Suppose two potentially
supersingular elliptic curves $E,E'/K$ have isomorphic nine-torsion,
common nontrivial unipotent residual image $C_3$, and different
prime-to-three special-inertia characters. Then their common
three-division field is generated by a cube root of a noncube
\emph{unit} of $K$. Their tame special-inertia quotient orders are
two and four, and their $j$-valuations, in some order, are
\[
                            3\quad\hbox{and}\quad b\ge6,
                                                        \tag{B-RD39.2}
\]
where $b=+\infty$ is allowed.

Consequently every remaining level-$27$ comparison with this residual
image has individual faithful special-inertia orders $12$ and $6$,
joint faithful order $36$, and independent wild order-three images.
Combined with Proposition~\ref{prop:B-RD38}, after the simultaneous
quadratic normalization of Proposition~\ref{prop:B-RD37}, the entire
remaining level-$27$ test has the following necessary alternatives:
\[
\begin{array}{c|c|c|c}
\text{residual image}&\text{individual special orders}&
 \text{joint order}&\text{$j$-valuations}\\\hline
C_4&(3,6)&18&(6,6)\\
C_3&(12,6)&36&(3,b),\ b\ge6.
\end{array}                                             \tag{B-RD39.3}
\]
The table asserts necessary restrictions, not existence of either kind
of comparison. Undoing the common quadratic twist may interchange
the tame orders one and two; it leaves order four and both
$j$-invariants unchanged.
\end{proposition}
\begin{proof}
A nontrivial unipotent residual image fixes a nonzero three-torsion
point. By \citet[Lemmas~4.1--4.2, p.~8]{Fisher2015Nine}, write
\[
 E=E_{t_E,D}:\quad y^2+3t_Exy+(t_E^3-D)y=x^3,
 \qquad L=K(E[3])=K(\sqrt[3]D).
                                                        \tag{B-RD39.4}
\]
Concretely, in the rational-three-point form
$y^2+a_1xy+a_3y=x^3$, take $t_E=a_1/3$ and
$D=t_E^3-a_3$. Its actual discriminant is $(3a_3)^3D$,
so its cube class is precisely the parameter class $D$.
The element $D$ is not a cube. A Weierstrass scaling replaces
$(t_E,D)$ by $(t_E/c,D/c^3)$, so we normalize
$d=v_K(D)\in\{0,1,2\}$. The common residual field $L/K$ is
cyclic and totally ramified of degree three. Fix $\theta\in L$
with $\theta^3=D$; then $v_L(\theta)=d$ and
$v_L(\zeta-1)=3$.

An arbitrary nine-congruence can be made direct or reverse by a scalar
change, as in Proposition~\ref{prop:B-RD37}. Fisher's exact
Theorem~4.3, equations~(10)--(11), pp.~9--10, supplies a parameter
$t\in K$ for $E'$ and requires
\[
\begin{array}{ll}
 \displaystyle\frac{t-\theta}{t_E-\theta}\in(L^\times)^3,
   &E'=E_{t,D}\quad\text{(direct)},\\[6pt]
 \displaystyle\frac{(t-\theta^{-1})(t_E-\theta)}3
                 \in(L^\times)^3,
   &E'=E_{t,D^{-1}}\quad\text{(reverse)}.
\end{array}                                             \tag{B-RD39.5}
\]
Each condition also holds with $\theta$ replaced by $\zeta\theta$
or $\zeta^2\theta$. The reverse denominator three is retained.

Here are all the required valuations, with no minimal-model
assumption. For $r=v_K(t)$ and $S=v_K(t^3-D)$, the exact invariants are
\[
 c_4(E_{t,D})=9t(t^3+8D),\qquad
 \Delta(E_{t,D})=27D(t^3-D)^3.
                                                        \tag{B-RD39.6}
\]
If $d=1$ or $2$, the valuations $3r$ and $d$ are unequal.
Potential supersingularity requires $r\ge-1$, and
\[
             S=\min(3r,d),\quad
             v_K(j)=6+3r-d,\quad
             \delta=6+d+3S.
                                                        \tag{B-RD39.7}
\]
The value $t=0$ belongs to the last row below and has $j=0$.
By Proposition~\ref{prop:B-RD36}, the tame quotient is read from
$\delta$ modulo four. The full direct-source tables are
\[
\begin{array}{c|ccc|ccc}
 &\multicolumn{3}{c|}{d=1}&\multicolumn{3}{c}{d=2}\\
r&S&\delta&\text{tame}&S&\delta&\text{tame}\\\hline
-1&-3&-2&2&-3&-1&4\\
0&0&7&4&0&8&1\\
\ge1&1&10&2&2&14&2
\end{array}                                             \tag{B-RD39.8}
\]
For the reverse target, whose parameter is $D^{-1}$, put
$S'=v_K(t^3-D^{-1})$. The analogous complete table is
\[
\begin{array}{c|ccc|ccc}
 &\multicolumn{3}{c|}{d=1}&\multicolumn{3}{c}{d=2}\\
r&S'&\delta'&\text{tame}&S'&\delta'&\text{tame}\\\hline
-2&-6&-13&4&-6&-14&2\\
-1&-3&-4&1&-3&-5&4\\
\ge0&-1&2&2&-2&-2&2
\end{array}                                             \tag{B-RD39.9}
\]

Take norms in (B-RD39.5). Since a norm of a cube is a cube in $K$,
the necessary congruences are
\[
    S-S_E\equiv0\pmod3\quad\text{(direct)},\qquad
    S'+S_E\equiv0\pmod3\quad\text{(reverse)}.
                                                        \tag{B-RD39.10}
\]
In the reverse case the norm of the denominator is $27$, whose
$K$-valuation is six, hence zero modulo three.

Suppose first $d=1$ or $2$. In the direct case, the only
opposite-tame pairs compatible with (B-RD39.10) have
$\{S,S_E\}=\{0,-3\}$. Interchange the pair if needed so that
$S_E=-3$, $S=0$. The quotient of the cube conditions at
$\zeta\theta$ and $\theta$ is
\[
 \frac{1+(1-\zeta)\theta/(t-\theta)}
      {1+(1-\zeta)\theta/(t_E-\theta)}.
                                                        \tag{B-RD39.11}
\]
The two displayed deviations have $L$-valuations $3+d$ and $6+d$.
The quotient is congruent to one with deviation valuation exactly
$3+d$, which is four or five. It cannot be a cube by
Lemma~\ref{lem:B-RD39-units}.

In the reverse case, compare the two tables using (B-RD39.10).
All opposite-tame pairs and the two deviation valuations in their
cube quotient are as follows:
\[
\begin{array}{c|rr|rr}
d&S_E&S'&3-d-S'&3+d-S_E\\\hline
1&-3&-6&8&7\\
1&-3&-3&5&7\\
1&0&-3&5&4\\
2&-3&-6&7&8\\
2&0&-6&7&5\\
2&0&-3&4&5
\end{array}                                             \tag{B-RD39.12}
\]
Indeed, on dividing the condition at $\zeta\theta$ by that at
$\theta$, one obtains the product
\[
 \left(1+\frac{(1-\zeta^{-1})\theta^{-1}}
                      {t-\theta^{-1}}\right)
 \left(1+\frac{(1-\zeta)\theta}{t_E-\theta}\right).
                                                        \tag{B-RD39.13}
\]
The valuations of its two deviations are the last two columns.
They are unequal; their minimum is four, five, or seven. The cross
term has strictly larger valuation. Thus the deviation of this
required cube from one has valuation four, five, or seven, again
contrary to Lemma~\ref{lem:B-RD39-units}. This excludes $d=1,2$
for either orientation.

We are left with the noncube unit $D$, of depth $s=1$ or $2$ in
Lemma~\ref{lem:B-RD39-units}. Its inverse has the same depth.
For either of the parameters $D,D^{-1}$, potential
supersingularity and the unit-depth lemma give exactly the following
possibilities:
\[
\begin{array}{c|c|c|c}
S&\delta&\text{tame}&v_K(j)\\\hline
-3&-3&4&3\\
0&6&2&\ge6\\
1\ (s=1)&9&4&6\\
2\ (s=2)&12&1&6
\end{array}                                             \tag{B-RD39.14}
\]
For $S=-3$ one has $v_K(t)=-1$. For $S=0$ the parameter is
integral and its residue does not equal the cube root of the unit
parameter; the zero value of $t$ is allowed. If $S>0$, $t$ is a
unit and $v_K(t^3+8D)=S$, since $S< v_K(9D)=4$.
These observations and (B-RD39.6) verify the table.

The direct norm congruence in (B-RD39.10) allows different tame
columns only for $\{S,S_E\}=\{-3,0\}$. The reverse norm
congruence likewise allows different columns only for that pair:
both curves have the same positive depth $s$, and $2s$ is not
divisible by three. Formula (B-RD39.14) now proves (B-RD39.2) and
the tame orders two and four.

At level $27$, Theorem~\ref{thm:B-RD34} restricts the joint faithful
order to eighteen or thirty-six. A tame quotient of order four
rules out eighteen, so the joint order is thirty-six. Its wild
subgroup has order nine; the two individual wild quotients have
order three and are independent. Their distinct tame orders four
and two therefore give individual orders twelve and six. The
other row of (B-RD39.3) is Proposition~\ref{prop:B-RD38}.
\end{proof}

\paragraph{Exact certificate and scope.}
The program \path{code/verify_c3_residual_gate.py} checks the
Weierstrass invariants, the exact Hensel polynomial, both source
and reverse-target tables, every opposite-tame norm-congruence
case, and all the required forbidden cube deviations. Its output
records the author runtime $0.053551534001599066$ seconds. Independent
author-output-preserving reruns by the root and second reviewers passed
in $0.045766205999825615$ and $0.05912446099682711$ seconds.
These calculations narrow the remaining test;
they do not construct or exclude either final row of (B-RD39.3).
